\chapter{Tipologie di azienda in Italia}

\begin{redbar}
Secondo l'articolo 2555 del Codice Civile
\begin{center}
\textit{"L’\textbf{azienda}\index{Organizzazione aziendale!azienda} è il complesso dei beni organizzati dall’imprenditore per l’esercizio dell’impresa"}.    
\end{center}
\end{redbar}

In Italia, le aziende si distinguono in due categorie principali:
\begin{enumerate}
    \item \textit{Aziende di produzione (imprese profit):}\index{Organizzazione aziendale!imprese profit} Queste aziende mirano a soddisfare i bisogni dei consumatori indirettamente, offrendo prodotti o servizi scambiati sul mercato. Si concentrano sul profitto e sull'efficienza nella produzione e distribuzione.
    \item \textit{Aziende di erogazione (corporazioni e fondazioni):} Operano per soddisfare direttamente i bisogni di determinati gruppi, come associati o fondatori, senza un obiettivo di profitto. Nelle \textit{corporazioni}\index{Organizzazione aziendale!corporazioni} prevale l’elemento personale e sono organizzate come associazioni, dove gli associati finanziano l’attività con contributi per il proprio beneficio. Nelle \textit{fondazioni}\index{Organizzazione aziendale!fondazioni}, invece, l’aspetto patrimoniale è centrale, e si sviluppano grazie a donazioni o lasciti destinati a obiettivi di utilità sociale e solidaristica.
\end{enumerate}

\section{Aziende individuali e società}
Le aziende di produzione possono essere gestite in forma individuale o collettiva. Nel caso di una \textit{ditta individuale}\index{Organizzazione aziendale!ditta individuale}, il proprietario è responsabile in modo illimitato per i debiti dell'azienda, senza distinzione tra il patrimonio personale e quello aziendale. Questo comporta vantaggi in termini di semplicità fiscale e operativa, ma anche rischi elevati a causa della responsabilità totale.

\begin{redbar}
Secondo l'articolo 2247 del Codice Civile
\begin{center}
\textit{"Con il contratto di \textbf{società}\index{Organizzazione aziendale!Società} due o più persone conferiscono beni o servizi per l’esercizio comune di un’attività economica allo scopo di dividerne gli utili"}.    
\end{center}
\end{redbar}

Le società si dividono in due tipi:
\begin{enumerate}
    \item \textit{Società di persone:}\index{Organizzazione aziendale!Società!persone} Comprendono la società semplice (s.s.), la società in nome collettivo (s.n.c.) e la società in accomandita semplice (s.a.s.). I soci sono responsabili illimitatamente per le obbligazioni sociali, ma il loro coinvolgimento diretto è spesso maggiore rispetto a quello nelle società di capitali.
    \item \textit{Società di capitali:}\index{Organizzazione aziendale!Società!capitali} Comprendono la società per azioni (s.p.a.), la società in accomandita per azioni (s.a.p.a.), la società a responsabilità limitata (s.r.l.) e la s.r.l. semplificata (s.r.l.s.). I soci non rispondono personalmente per le obbligazioni sociali, limitando la responsabilità al capitale conferito. Questo tipo di società facilita l’attrazione di capitali esterni ma è soggetto a una doppia imposizione fiscale e a una gestione più onerosa.
\end{enumerate}

\section{Soggetto giuridico e soggetto economico}
Nelle aziende, il \textit{soggetto giuridico}\index{Organizzazione aziendale!soggetto giuridico} è l'entità che assume diritti e obblighi legali derivanti dall’attività aziendale. Può essere una persona fisica, come nel caso di una ditta individuale, o una società, che gode di una propria personalità giuridica e capacità di essere titolare di diritti distinti rispetto a quelli dei soci.

Il \textit{soggetto economico}\index{Organizzazione aziendale!soggetto economico}, invece, è chi effettivamente controlla e dirige l'azienda. Nelle ditte individuali, è il proprietario; nelle società di persone, i soci; nelle società di capitali, può essere la maggioranza del capitale sociale o, in aziende molto grandi (public company), i manager stessi.

\section{Elementi costitutivi dell’azienda}
Il capitale finanziario di un’azienda si suddivide in:
\begin{enumerate}
    \item \textit{Capitale fisso o immobilizzato:}\index{Organizzazione aziendale!Capitale!fisso o immobilizzato} Include immobilizzazioni materiali (come fabbricati e attrezzature), immobilizzazioni immateriali (come brevetti e marchi) e immobilizzazioni finanziarie (come partecipazioni e titoli).
    \item \textit{Capitale circolante:}\index{Organizzazione aziendale!Capitale!cicolante} Comprende il magazzino, i crediti verso clienti e le disponibilità liquide, come cassa e depositi bancari.
\end{enumerate}
Insieme, capitale fisso e circolante costituiscono il \textit{capitale di funzionamento}\index{Organizzazione aziendale!Capitale!funzionamento} dell’impresa (c.d.), essenziale per la sua operatività quotidiana.